Long-term experiments must choose a costly subset of short-term outcomes, attach sequential proxy roles, and allocate a fixed main-study budget across experimental short-term units and observational long-term follow-up. In the persistent-confounding model of Imbens, Kallus, Mao, and Wang, this note makes that measurement design the argument of a two-source score-variance profile. Its named negative result transports the Carneiro--Lee--Wilhelm balanced-treatment population criterion under the exact linear paired-unit cost specialization: using a nonredundant affine basis and a nonsingular full Gram matrix, the rule minimizes affine-linear residual variance per budget. An explicit unequal-cost binary family proves that this named population rule can select a design whose arm-normalized and IKM-normalized bridge values have an unbounded ratio to the optimum on a relatively open regular, common-identifying region; the exact integer floor-budget selector agrees for every fixed law once the budget exceeds a finite threshold. Generic unrestricted unadjusted and treatment-aware prediction-risk failures remain separate secondary results. A full-catalogue pilot can instead enumerate estimated source-score variances, select the minimum bridge design value, implement the square-root source allocation, and support selected-design Wald inference under the stated learner and variance conditions.